%==============================================================================
% TRACE-MAS v3: Securing the Inter-Agent Surface in Multi-Agent LLM Systems
% Author: S. Gentyala
% Format: IEEEtran (conference, two-column)
% Compile: pdflatex -> pdflatex
% GitHub: https://github.com/sunilgentyala/TRACE-MAS
%==============================================================================
\documentclass[conference]{IEEEtran}
\IEEEoverridecommandlockouts

\usepackage{cite}
\usepackage{amsmath,amssymb,amsfonts,amsthm}
\usepackage{algorithmic}
\usepackage{algorithm}
\usepackage{graphicx}
\usepackage{textcomp}
\usepackage{xcolor}
\usepackage{url}
\usepackage{hyperref}

\hypersetup{
  colorlinks=true,
  linkcolor=blue,
  citecolor=blue,
  urlcolor=blue
}

\newtheorem{theorem}{Theorem}
\newtheorem{lemma}{Lemma}
\newtheorem{definition}{Definition}
\newtheorem{proposition}{Proposition}
\newtheorem{assumption}{Assumption}
\newtheorem{remark}{Remark}

\begin{document}

\title{TRACE-MAS: Securing the Inter-Agent Surface Against Cascading Drift,
Identity Spoofing, Prompt Injection, and Temporal Adversarial Attacks
in Multi-Agent LLM Pipelines}

\author{%
\IEEEauthorblockN{Sunil Gentyala}
\IEEEauthorblockA{%
\textit{Independent Research, MAS Security and Applied Cryptography}\\
Email: sugentyala@ieee.org\\
GitHub: \url{https://github.com/sunilgentyala/TRACE-MAS}}
}

\maketitle

\begin{abstract}
Modern AI systems rarely rely on a single model. A coder agent
drafts a module, a quality-assurance agent reviews it, and a deployment
agent ships it---three distinct LLMs sharing a common pipeline but
operating without any formal cross-verification layer. When that pipeline
is attacked, the consequences are not isolated to one model; they propagate
invisibly across the entire chain. This paper identifies four structurally
related failure modes that define the attack surface of such pipelines:
\emph{cascading alignment drift}, \emph{inter-agent identity spoofing},
\emph{cross-agent prompt injection}, and \emph{temporal adversarial
manipulation}. To address all four jointly, we introduce
\textsc{TRACE-MAS} (\emph{T}ri-vector \emph{R}esilient
\emph{A}lgorithm for \emph{C}ooperative \emph{E}mbodied Multi-Agent
Security), a unified theoretical framework grounded in contraction theory,
zero-knowledge cryptography, and statistical change detection.

\textsc{TRACE-MAS} delivers three provable guarantees. First, a
verifier-mediated contractive aggregation operator bounds the asymptotic
drift of any $n$-agent chain to $\Delta_\infty \leq \rho + \lambda\epsilon/(1-\lambda)$,
independent of chain length. Second, a dynamic zero-knowledge verification
threshold $\tau_i^{(t)}$ that adapts to the runtime risk of each agent
bounds the probability that a rogue participant is accepted in any round to
$2^{-s} + e^{-\tau/\mu}$; the same mechanism rejects messages whose
content diverges from the agent's declared policy commitment, which is
the cryptographic footprint of a prompt injection event. Third, a
sliding-window Kullback--Leibler functional over joint agent memory
detects slow temporal attacks with a sub-Gaussian miss probability
$\exp(-T_w(\varphi(B) - \nu)^2 / 2\sigma^2)$.
Together, the three theorems compose into a single security envelope
whose failure terms are independently controllable by parameter choice.
We present full algorithmic pseudocode, an asymptotic cost analysis,
and a discussion of open engineering challenges in applying the
framework to production agentic pipelines.
\end{abstract}

\begin{IEEEkeywords}
multi-agent systems, alignment drift, zero-knowledge proofs, prompt
injection, agent-to-agent authentication, temporal adversarial attacks,
contraction theory, LLM pipeline security.
\end{IEEEkeywords}

%=============================================================================
\section{Introduction}
\label{sec:intro}
%=============================================================================

The shift from single-model AI to cooperative multi-agent pipelines has
been one of the defining architectural transitions of the past two years.
Standardized agent communication protocols such as Google's Agent2Agent
(A2A), Anthropic's Model Context Protocol, and OpenAI's Swarm framework
all reflect the same
engineering judgement: complex tasks are solved better by specialized
agents coordinating over a shared interface than by a single monolithic
model trying to do everything. Production deployments now chain together
coder agents, retrieval agents, QA agents, deployment agents, and
financial-transaction agents in pipelines that operate with minimal human
oversight. The throughput is impressive. The attack surface is alarming.

The inter-agent surface---the seams between those specialized agents---is
the least defended layer of the agentic stack. An agent boundary is not a
security boundary. When agent $A_1$ hands its output to agent $A_2$, it
hands over whatever drift, corruption, or injected instruction that output
carries. $A_2$, lacking the original intent, treats the input as ground
truth. $A_3$ does the same for $A_2$'s output. This propagation dynamic
means that a small attack at the front of a pipeline can produce a large
effect at its end---a property we call \emph{cascading alignment drift}.

Three additional failure modes compound the problem. Platforms that allow
open-world agent participation---where any service advertising the right
capability descriptor can join a workflow---face \emph{inter-agent
identity spoofing}: a rogue agent is admitted, issues commands that
downstream agents execute in good faith, and exits before anomaly
detection fires. Separately, the cross-agent analogue of prompt injection
is now well-documented~\cite{ref:promptinfection2024}: an adversary
embeds a directive in content processed by $A_1$ that survives intact
into $A_2$'s context, effectively hijacking a downstream agent's
reasoning without compromising the upstream one. Finally, patient
adversaries can engineer \emph{temporal adversarial attacks} that stage
incremental perturbations whose individual magnitudes remain below any
single agent's anomaly threshold but whose cumulative effect---integrated
across the agent's memory horizon---drives the system toward a prepared
malicious end-state~\cite{ref:agentpoison2024}.

Existing defenses are compartmentalized. Alignment research tackles drift
in isolation. Cryptographic authentication schemes address identity but
not intent. Prompt-injection defenses operate at the input layer of a
single model without awareness of multi-agent propagation dynamics.
Temporal robustness work in adversarial ML handles single-agent sequential
decision settings. None of these threads considers all four attack vectors
jointly, and it is precisely their interaction---drift enabling injection,
injection enabling spoofing escalation, and staged perturbations exploiting
the window blindness of each individual defense---that makes the
inter-agent surface so dangerous.

We close that integration gap. Our contribution is \textsc{TRACE-MAS}
(\emph{T}ri-vector \emph{R}esilient \emph{A}lgorithm for \emph{C}ooperative
\emph{E}mbodied Multi-Agent Security), a unified theoretical framework that
addresses all four failure modes under one composable security envelope.
The framework's three theorems cover (i) drift decay under contractive
verifier-mediated aggregation, (ii) spoofing-and-injection acceptance
bounds via dynamic zero-knowledge verification thresholds, and
(iii) temporal attack detection via a sliding-window divergence functional.
The composite failure probability is proven to be independently controllable
on each dimension.

The remainder of this paper is organized as follows.
Section~\ref{sec:related} surveys related work on multi-agent reliability,
cryptographic authentication, prompt injection, and temporal adversarial
attacks, and isolates the integration gap.
Section~\ref{sec:framework} develops the mathematical framework:
system model, the three theorems with proof sketches, and the new subsection
on LLM-specific injection vectors.
Section~\ref{sec:algo} presents the unified algorithm.
Section~\ref{sec:discussion} examines verifier collusion, calibration
drift, embodied extensions, privacy, and limitations.
Section~\ref{sec:conclusion} summarizes and outlines a research roadmap.

%=============================================================================
\section{Related Work and Threat Landscape}
\label{sec:related}
%=============================================================================

\subsection{Multi-Agent Reliability and Consensus}
The recent review of embodied multi-agent systems~\cite{ref:embodied2025}
consolidates a decade of cooperative-control work and identifies
cross-agent verification as the most pressing open frontier in the
field. Resilient consensus tracking under semi-Markovian switching
topologies~\cite{ref:jas2025consensus} provides one mathematical
foundation, but assumes correctness of agent identity---an assumption
this paper explicitly drops. Byzantine fault-tolerant consensus
protocols~\cite{ref:tc2024pbft} supply the quorum primitive, yet are
calibrated for closed permissioned networks rather than the capability-
discovered swarms that characterize modern agentic deployments. The
privacy-distributed optimization work of~\cite{ref:jas2024priv}
protects data flows but leaves the intent layer unguarded.

\subsection{Adversarial Attacks in Multi-Agent Settings}
Grey-box adversarial attacks on communicative multi-agent reinforcement
learning demonstrate that partial knowledge of policy structure is
sufficient to craft message-level subversions~\cite{ref:tifs2025grey}.
Safe multi-agent RL under adversarial communications offers a wireless
defense but presupposes agent provenance~\cite{ref:tifs2024safe}. The
contemporary survey of adversarial attacks in ML-empowered communication
systems~\cite{ref:comst2023} frames the sequential perturbation regime
that motivates our temporal-window analysis. None of these threads
addresses the compound attack surface that emerges when drift, spoofing,
injection, and temporal manipulation occur concurrently.

\subsection{Prompt Injection and Memory Poisoning}
A qualitatively new class of attack has emerged as LLM agents entered
production. \emph{Prompt infection}~\cite{ref:promptinfection2024}
demonstrates that a single injected directive in content processed by
one agent propagates virally to downstream agents in the same pipeline,
coordinating them to exfiltrate data or issue unauthorized commands.
\emph{AgentPoison}~\cite{ref:agentpoison2024} attacks the retrieval
layer: by poisoning a small fraction of an agent's long-term memory or
RAG knowledge base, an adversary can redirect behavior with near-perfect
reliability and minimal impact on benign queries. DataSentinel
addresses this for single-agent settings via game-theoretic
minimax detection~\cite{ref:datasentinel2025}, but a multi-agent
extension does not yet exist. Section~\ref{subsec:injection} shows
that \textsc{TRACE-MAS}'s ZKP attestation provides exactly that
extension: a message whose content violates the sending agent's
declared policy commitment fails verification by construction.

\subsection{Cryptographic Authentication for Agents}
Zero-knowledge mutual authentication schemes engineered for IoT
networks~\cite{ref:zpma2024} provide a cryptographic backbone, but
their identity-only proofs do not cover behavioral commitments.
Privacy-preserving tool-use harnesses for LLM agents via homomorphic
encryption~\cite{ref:tdsc2024} protect data without guarding intent.
Secure agent communication protocols are an active standards concern
in 2025~\cite{ref:agentsec2025}, with proposals converging on
OAuth-anchored attestation but lacking provable behavioral soundness.
\textsc{TRACE-MAS}'s dynamic ZKP threshold generalizes these
approaches to message-level policy commitments.

\subsection{Multi-Agent Systems for Security Applications}
The inverse direction---agents acting as defenders rather than as
attack surfaces---is documented in~\cite{ref:access2025multi}, which
applies MAS to cybersecurity threat detection and correlation. That
work motivates why securing the agents themselves is critical: a
multi-agent defense system whose inter-agent surface is compromised
is worse than no defense at all.

%=============================================================================
\section{Mathematical Framework}
\label{sec:framework}
%=============================================================================

\subsection{System Model and Notation}
\label{subsec:model}

Let $\mathcal{A} = \{A_1, A_2, \ldots, A_n\}$ denote a chain or swarm of
agents executing a workflow $\mathcal{W}$. At round
$t \in \{1, \ldots, T\}$, agent $A_i$ produces an output
$h_i^{(t)} \in \mathbb{R}^d$ in a shared embedding space. Let
$\alpha_i^{(t)}$ denote the canonical aligned representation that
$A_i$ should produce given the user intent $\alpha_0$. Let $M^{(t)}$
denote the joint memory of the swarm, and let
$\mathcal{V} = \{V_1, \ldots, V_k\}$ denote a quorum of independent
verifier agents producing a reference signal $\xi^{(t)}$.

\begin{definition}[Inter-Agent Surface]
The inter-agent surface $\mathcal{S}$ is the tuple
\[
\mathcal{S} = \bigl(\{h_i^{(t)}\},\ \{\pi_i^{(t)}\},\
M^{(t)},\ \xi^{(t)}\bigr),
\]
where $\pi_i^{(t)}$ is the cryptographic attestation accompanying
$h_i^{(t)}$.
\end{definition}

The three core security gaps target three components of $\mathcal{S}$:
drift attacks corrupt $\{h_i^{(t)}\}$; spoofing attacks forge or omit
$\{\pi_i^{(t)}\}$; temporal attacks manipulate $M^{(t)}$ over horizons
longer than any individual agent's anomaly window. Cross-agent prompt
injection is a specialization of spoofing: the forged element is not
identity but intent, and the attestation mechanism of
Section~\ref{subsec:zkp} is the mechanism that catches it.

\subsection{Drift-Decay Bounding Function}
\label{subsec:drift}

We model alignment drift at handoff $i$ as
$\delta_i \triangleq \|\alpha_i - \alpha_0\|_2$. In an unprotected
chain, each agent injects a per-step error $\varepsilon_i$ with
$\|\varepsilon_i\|_2 \leq \epsilon$, yielding the loose bound
$\delta_n \leq \delta_0 + n\epsilon$ that grows without bound in
chain length. We replace pointwise propagation with a verifier-mediated
contractive aggregation operator $\mathcal{T}$.

\begin{definition}[Verifier-Mediated Aggregation]
Let $\gamma_i \in [\gamma_{\min}, 1)$ be the verifier weight at step
$i$, and let the verifier signal $\xi_i$ satisfy
$\|\xi_i - \alpha_0\|_2 \leq \rho$. The aggregation operator is
\begin{equation}
\alpha_{i+1} = (1-\gamma_i)\hat{\alpha}_{i+1} + \gamma_i \xi_i,
\end{equation}
where $\hat{\alpha}_{i+1}$ is the natural agent output satisfying
$\|\hat{\alpha}_{i+1} - \alpha_i\|_2 \leq \epsilon$.
\end{definition}

\begin{theorem}[Drift-Decay Theorem]
Under Definition~2 with $\lambda \triangleq 1 - \gamma_{\min} \in (0,1)$,
the drift sequence satisfies
\begin{equation}
\delta_{i+1} \leq \lambda\,\delta_i + \lambda\epsilon +
(1-\lambda)\rho
\end{equation}
and consequently
\begin{equation}
\delta_n \leq \lambda^n \delta_0 +
\frac{\lambda\epsilon + (1-\lambda)\rho}{1-\lambda}
\bigl(1-\lambda^n\bigr).
\end{equation}
In particular, the asymptotic drift envelope is
\begin{equation}
\delta_\infty \leq \rho + \frac{\lambda\epsilon}{1-\lambda},
\end{equation}
which is \emph{independent of chain length $n$}.
\end{theorem}

\begin{IEEEproof}[Proof Sketch]
Apply the triangle inequality to (1):
\begin{align*}
\delta_{i+1} &= \|\alpha_{i+1} - \alpha_0\|
\leq (1-\gamma_i)\|\hat{\alpha}_{i+1} - \alpha_0\| +
\gamma_i\|\xi_i - \alpha_0\|\\
&\leq (1-\gamma_i)(\delta_i + \epsilon) + \gamma_i \rho
\leq \lambda(\delta_i + \epsilon) + (1-\lambda)\rho.
\end{align*}
Iterating and summing the geometric series yields~(3).
Taking $n \to \infty$ delivers~(4).
\end{IEEEproof}

The practical consequence is significant: the cascading-drift
catastrophe familiar from unverified LLM chains is replaced by a
fixed budget that depends on verifier quality $\rho$ and per-agent
noise $\epsilon$, but not on how many agents the workflow chains
together. This means that longer pipelines are not inherently riskier
under \textsc{TRACE-MAS}, provided the verifier quorum quality
is maintained.

\subsection{Dynamic ZKP Verification Threshold}
\label{subsec:zkp}

To neutralize inter-agent spoofing, every message-bearing handoff must
carry a non-interactive zero-knowledge proof that jointly binds identity,
policy commitment, and message provenance.

\begin{definition}[Agent Attestation]
Let $sk_i, pk_i$ denote agent $A_i$'s signing key pair, and let
$C_{\theta_i}$ denote a hiding commitment to $A_i$'s declared policy
parameters $\theta_i$. The agent generates
\[
\pi_i^{(t)} \leftarrow \mathsf{Prove}\bigl(sk_i,\, m_i^{(t)},\,
C_{\theta_i},\, \mathcal{H}_i^{(t)}\bigr),
\]
attesting jointly to: (i) knowledge of $sk_i$,
(ii) consistency of message $m_i^{(t)}$ with declared policy $\theta_i$,
and (iii) consistency of $\theta_i$ with the agent's observed history
$\mathcal{H}_i^{(t)}$.
\end{definition}

The verifier accepts $\pi_i^{(t)}$ iff the standard ZKP verification
passes and the behavioral score crosses a dynamic threshold:
\begin{equation}
\mathsf{Verify}(\pi_i^{(t)}) = 1 \ \wedge\
\mathsf{Score}\bigl(\pi_i^{(t)}, \mathcal{H}_i^{(t)}\bigr) \geq
\tau_i^{(t)}.
\end{equation}

\begin{definition}[Dynamic Threshold]
Let $\rho_i^{(t)} \in [0, \infty)$ be the runtime risk score of agent
$i$ at time $t$, computed from prior anomalies, blast radius of the
pending action, and network exposure. The dynamic threshold is
\begin{equation}
\tau_i^{(t)} = \tau_0 + \kappa\,\log\bigl(1 + \rho_i^{(t)}\bigr).
\end{equation}
\end{definition}

\begin{theorem}[Spoofing-Acceptance Bound]
Let the underlying ZKP scheme have soundness error $2^{-s}$, and
let an adversary forging behavioral evidence enjoy advantage at most
$\mu > 0$ in producing high-scoring transcripts. For any single
round,
\begin{equation}
\Pr[\text{rogue agent accepted}] \leq 2^{-s} + e^{-\tau_i^{(t)}/\mu}.
\end{equation}
Over $N$ rounds, with the threshold calibrated as
$\tau_i^{(t)} \geq \mu\log(1/\varepsilon)$, the swarm's
union-bounded spoofing probability is at most
$N(2^{-s} + \varepsilon)$.
\end{theorem}

\begin{IEEEproof}[Proof Sketch]
The first term follows directly from soundness of the ZKP scheme.
The second follows from a sub-exponential tail bound on $\mathsf{Score}$
under the assumption that the adversary's transcript distribution
satisfies $\Pr[\mathsf{Score} \geq \tau] \leq e^{-\tau/\mu}$.
The union bound across $N$ rounds completes the argument.
\end{IEEEproof}

The adaptive shape of $\tau_i^{(t)}$ is central to deployability.
Routine, low-risk handoffs---passing intermediate results between
retrieval and reasoning agents---incur minimal verification overhead.
High-risk operations---payments, deployments, privilege escalations---
automatically demand cryptographic proofs of correspondingly higher
confidence. This inherits from blockchain-anchored mutual-authentication
schemes~\cite{ref:zpma2024} but generalizes them to behavioral
commitments rather than mere identity.

\subsection{Temporal Memory-Validation Boundary}
\label{subsec:temporal}

Temporal attacks succeed when no agent's observation window is long
enough to detect the cumulative perturbation. We therefore impose a
sliding-window inconsistency functional on the joint memory.

\begin{definition}[Temporal Inconsistency Functional]
For window size $T_w$ centered at time $t$, define
\begin{equation}
I(t, T_w) \triangleq \sup_{(s_1, s_2) \in \mathcal{W}(t,T_w)}
D_{\mathrm{KL}}\bigl(M^{(s_1)} \,\big\|\, M^{(s_2)}\bigr),
\end{equation}
where $\mathcal{W}(t, T_w)$ is the set of admissible time pairs in
the window.
\end{definition}

\begin{definition}[Temporal Validation Boundary]
Given an adversarial perturbation budget $\|\eta\|_p \leq B$ and a
calibration function $\varphi(\cdot)$, the validation boundary is
the optimal window length $T_w^{*}$ satisfying
\begin{equation}
I\bigl(t, T_w^{*}\bigr) \leq \varphi(B)
\ \Leftrightarrow\ \text{benign trajectory}.
\end{equation}
Any window in which $I(t, T_w^{*}) > \varphi(B)$ triggers temporal
quarantine.
\end{definition}

\begin{assumption}[Sub-Gaussian Benign Drift]
Benign memory increments are sub-Gaussian with mean shift $\nu$ and
variance proxy $\sigma^2$.
\end{assumption}

\begin{theorem}[Temporal Detection Bound]
Under Assumption~1, the probability that an adversarial trajectory
with cumulative budget $B$ evades detection inside a window of length
$T_w$ satisfies
\begin{equation}
\Pr[\text{undetected at time } t] \leq
\exp\!\left(-\frac{T_w(\varphi(B) - \nu)^2}{2\sigma^2}\right).
\end{equation}
The optimal window solves
\begin{equation}
T_w^{*} = \arg\min_{T_w}\,\bigl\{\mathcal{R}_{\mathrm{FN}}(T_w, B)
+ \mathcal{R}_{\mathrm{FP}}(T_w)\bigr\}.
\end{equation}
\end{theorem}

\begin{IEEEproof}[Proof Sketch]
Apply Hoeffding's inequality to the sum of sub-Gaussian increments
under the null hypothesis of benign drift. The bound~(10) is a
standard concentration result. Equation~(11) follows from balancing
the false-negative/false-positive detection-error decomposition.
\end{IEEEproof}

\subsection{Resilience to LLM-Specific Injection Vectors}
\label{subsec:injection}

The three theorems above were stated for abstract agent outputs
$h_i^{(t)}$. LLM agents introduce a specific injection pathway that
does not arise in classical multi-agent control: an adversary can
embed a directive in \emph{content} that a legitimate agent processes
(a retrieved document, a tool output, an upstream message), causing
that agent to emit a message that serves the adversary's goal rather
than the declared workflow. This is the cross-agent analogue of prompt
injection and has been formalized as \emph{Prompt Infection}
by~\cite{ref:promptinfection2024}.

Under \textsc{TRACE-MAS}, this attack is covered by Theorem~2 without
modification. When a prompt injection event causes agent $A_i$ to
deviate from its declared policy $\theta_i$, the injected message
$m_i^{(t)}$ is inconsistent with $C_{\theta_i}$, and the attestation
proof $\pi_i^{(t)}$ cannot be generated for it---there is no witness.
The verifier therefore rejects the handoff.

\begin{remark}
The ZKP attestation in Definition~3 binds message content to a
policy commitment, not merely to an identity key. An agent that
generates a message under an adversarial directive cannot produce a
valid proof $\pi_i^{(t)}$ for that message unless the directive is
consistent with the agent's declared policy $\theta_i$. Prompt
injection attacks typically direct the agent to act inconsistently
with its declared role (e.g., exfiltrate data, issue unauthorized
commands), so the proof generation step fails by soundness.
\end{remark}

Memory poisoning attacks~\cite{ref:agentpoison2024} that corrupt
the retrieval layer are captured by Theorem~3: poisoned retrievals
introduce abnormal distributional shifts in $M^{(t)}$ that exceed
the benign-drift calibration $\varphi(B)$, triggering temporal
quarantine. The window length $T_w^{*}$ must be calibrated to the
attacker's poisoning rate, which is a standard estimation problem
(see limitations in Section~\ref{sec:discussion}).

\subsection{Composite Security Envelope}
\label{subsec:composite}

Combining Theorems~1, 2, and~3, the joint failure probability of
\textsc{TRACE-MAS} over a workflow of $n$ agents and $T$ rounds is
bounded by
\begin{equation}
P_{\mathrm{fail}} \leq
\mathbf{1}[\delta_\infty > \delta^{*}]
+ T\bigl(2^{-s} + \varepsilon\bigr)
+ T\,\exp\!\left(-\frac{T_w^{*}(\varphi(B) - \nu)^2}{2\sigma^2}
\right).
\end{equation}
Each term is independently controllable: the first by verifier quality
$\rho$ and contraction rate $\lambda$; the second by soundness
parameter $s$ and threshold calibration; the third by window length
$T_w^{*}$ and noise calibration $\sigma^2$. Cascading drift cannot
inflate the spoofing or temporal terms because the verifier-mediated
aggregation acts---and resets drift---before downstream propagation.

%=============================================================================
\section{Algorithmic Pseudocode}
\label{sec:algo}
%=============================================================================

The unified runtime executes the three protections in sequence per
round. The algorithm and all supporting materials are open-sourced at
\url{https://github.com/sunilgentyala/TRACE-MAS}.

\begin{algorithm}[t]
\caption{\textsc{TRACE-MAS} Unified Runtime}
\begin{algorithmic}[1]
\REQUIRE Agent chain $\mathcal{A}=\{A_1,\ldots,A_n\}$; verifier
quorum $\mathcal{V}$; user intent $\alpha_0$;
ZKP scheme $(\mathsf{Prove},\mathsf{Verify})$;
calibration $\varphi(\cdot)$; hyperparameters
$\gamma_{\min}, \tau_0, \kappa, T_w^{*}, B$
\ENSURE Certified trajectory $(\alpha_1,\ldots,\alpha_n)$ or
$\textsc{Halt}$
\STATE $\alpha_0 \leftarrow$ user intent;\;
$M^{(0)} \leftarrow \emptyset$
\FOR{$t = 1$ \TO $T$}
  \FOR{each agent $A_i$ acting at round $t$}
    \STATE $A_i$ produces $\hat{\alpha}_{i+1}$
    and message $m_i^{(t)}$
    \STATE $\pi_i^{(t)} \leftarrow \mathsf{Prove}(sk_i,
    m_i^{(t)}, C_{\theta_i}, \mathcal{H}_i^{(t)})$
    \STATE Compute $\rho_i^{(t)}$;\;
    $\tau_i^{(t)} \leftarrow \tau_0 +
    \kappa\log(1+\rho_i^{(t)})$
    \IF{$\mathsf{Verify}(\pi_i^{(t)}) \neq 1$ \OR
        $\mathsf{Score}(\pi_i^{(t)},\mathcal{H}_i^{(t)})
        < \tau_i^{(t)}$}
      \STATE Quarantine $A_i$;\;
      raise \textsc{SpoofOrInjectAlarm};\;
      \textbf{return} \textsc{Halt}
    \ENDIF
    \STATE $\mathcal{V}$ produces reference $\xi_i^{(t)}$
    \STATE $\gamma_i \leftarrow
    \max\bigl(\gamma_{\min},\,g(\rho_i^{(t)})\bigr)$
    \STATE $\alpha_{i+1} \leftarrow
    (1-\gamma_i)\hat{\alpha}_{i+1} + \gamma_i \xi_i^{(t)}$
    \COMMENT{Eq.~(1)}
    \STATE $\delta_{i+1} \leftarrow
    \|\alpha_{i+1} - \alpha_0\|_2$
    \IF{$\delta_{i+1} > \delta^{*}$}
      \STATE Raise \textsc{DriftAlarm};\;
      \textbf{return} \textsc{Halt}
    \ENDIF
    \STATE $M^{(t)} \leftarrow M^{(t-1)} \cup
    \{\alpha_{i+1}\}$
  \ENDFOR
  \STATE Compute $I(t, T_w^{*})$ via Eq.~(8)
  \IF{$I(t, T_w^{*}) > \varphi(B)$}
    \STATE Raise \textsc{TemporalAlarm};\;
    quarantine suspect window in $M$
    \STATE Roll back to last certified checkpoint;\;
    \textbf{return} \textsc{Halt}
  \ENDIF
\ENDFOR
\RETURN Certified trajectory $(\alpha_1,\ldots,\alpha_n)$
\end{algorithmic}
\end{algorithm}

\noindent\textbf{Complexity.}
Each round costs $O(\log|\mathcal{A}|)$ for ZKP verification
(succinct schemes), $O(d)$ for the aggregation step, and
$O(T_w^{*})$ for the window functional. The dominant overhead in
deployment is ZKP proof \emph{generation}, not verification;
this is amortizable through batched proof aggregation or
pre-committed SNARKs computed ahead of each workflow step.

%=============================================================================
\section{Discussion and Threat-Model Boundaries}
\label{sec:discussion}
%=============================================================================

\paragraph*{Verifier collusion}
Theorem~1 assumes the verifier residual $\rho$ is bounded. In
practice this requires either a quorum threshold
$k > k_{\mathrm{Byz}}$ with classical Byzantine
guarantees~\cite{ref:tc2024pbft} or a heterogeneous quorum sampled
from disjoint provider trust domains. Whether $k_{\mathrm{Byz}}$
remains meaningful in fully open swarms whose membership shifts on
the order of seconds is, in our view, an open question worth
dedicated study.

\paragraph*{Threshold calibration drift}
The dynamic threshold $\tau_i^{(t)}$ is itself an attack surface:
a patient adversary could nudge $\rho_i^{(t)}$ downward over time
to reduce the effective security bar. We mitigate this by requiring
the risk-score path itself to be temporally validated under
Theorem~3, closing the feedback loop. A complete fixed-point
characterization under mild monotonicity assumptions on $g(\cdot)$
is left to future work.

\paragraph*{Prompt injection and memory poisoning in production}
Recent attacks document that prompt injection in multi-agent
settings propagates with alarming efficiency~\cite{ref:promptinfection2024},
and memory poisoning can achieve high attack success rates at very
low poison fractions~\cite{ref:agentpoison2024}. The coverage
argument in Section~\ref{subsec:injection} shows that
\textsc{TRACE-MAS} intercepts both via the attestation layer.
However, the argument relies on the ZKP prover having access to the
policy commitment $C_{\theta_i}$ \emph{at proof time}. If an
adversary can corrupt $C_{\theta_i}$ before attestation begins---for
example, by poisoning the agent's initialization context---the
framework's injection guarantee does not apply. Protecting the
initialization phase is outside our current scope.

\paragraph*{Agent communication protocol alignment}
Emerging standards for agent-to-agent communication~\cite{ref:agentsec2025,ref:a2a2025}
converge on OAuth-anchored authentication over HTTPS, but do not yet
specify behavioral commitments or policy proofs. \textsc{TRACE-MAS}
is designed to be composable with these standards: the ZKP attestation
layer can be implemented as an additional header on A2A or MCP
messages, and the verifier quorum can run as a sidecar service that
intercepts inter-agent calls.

\paragraph*{Embodiment}
For embodied agents whose actions perturb the physical
environment~\cite{ref:embodied2025}, the temporal functional must
include observational state $o^{(t)}$ alongside memory state. The
framework extends without modification by replacing $M^{(t)}$ with
the joint $(M^{(t)}, o^{(t)})$, though a careful treatment of partial
observability deserves its own investigation.

\paragraph*{Privacy}
Privacy of agent reasoning is preserved by zero-knowledge: the
attestation reveals nothing about $\theta_i$ beyond its commitment.
Combined with prior work on tool-use privacy~\cite{ref:tdsc2024}
and federated optimization under privacy constraints~\cite{ref:jas2024priv},
\textsc{TRACE-MAS} is compatible with confidentiality requirements
in healthcare and finance.

\paragraph*{Limitations}
Several open challenges limit the immediate deployability of
\textsc{TRACE-MAS} as a production system. First, the framework
is theoretical: empirical validation across coder--QA--deploy
pipelines and embodied robotic swarms is reserved for follow-up
work, and the practical constants in our bounds will only become
trustworthy once measured against real workloads. Second, proof
generation at industrial LLM dialogue rates---which can exceed
hundreds of messages per second in large swarms---remains
computationally intensive, even for succinct ZKP
backends; batching and pre-commitment strategies reduce but do not
eliminate this overhead. Third, our verifier quorum model assumes
a fixed Byzantine budget $k_{\mathrm{Byz}}$; what happens when
the quorum itself is partially compromised or when membership
changes faster than the quorum can be reconstituted is an open
and under-explored question. Fourth, calibrating $\varphi(B)$
and $\sigma^2$ without ground-truth knowledge of the adversarial
budget requires online estimation, and doing so safely without
admitting calibration-time poisoning is nontrivial. Fifth,
Theorem~3 rests on sub-Gaussian benign drift; real-world
multi-day agent rollouts may exhibit heavier tails, and the bound's
sensitivity under those distributions deserves characterization.
Finally, the framework presumes correctness of the underlying
cryptographic primitives; cryptographic agility---the ability to
swap ZKP backends as advances in cryptanalysis emerge---is
desirable but not yet specified.

%=============================================================================
\section{Conclusion}
\label{sec:conclusion}
%=============================================================================

The security gaps in multi-agent LLM pipelines are not independent
vulnerabilities. They are a coherent failure family rooted in the
absence of verified cross-checking on the inter-agent surface. A
drift that goes undetected at step $i$ becomes the ground truth
for step $i+1$; an identity that goes unauthenticated admits rogue
participants; a policy that goes uncommitted admits prompt injection;
a memory that goes unmonitored admits slow temporal manipulation.
\textsc{TRACE-MAS} treats these four failure modes jointly and closes
them with a single composable framework.

The three core theorems deliver a chain-length-independent drift
envelope, a logarithmically-tight spoofing-and-injection acceptance
bound, and a sub-Gaussian temporal-detection bound. Their composition
yields a failure probability whose three terms are independently
controllable. The unified algorithm runs in $O(\log|\mathcal{A}|
+ d + T_w^{*})$ per round, with ZKP proof generation as the dominant
amortizable cost.

The next phase of this research program is empirical: instantiation
of the ZKP attestation layer using a concrete succinct proof system
(e.g., Groth16 or PLONK), calibration of $\varphi(B)$ on
representative benchmark pipelines, red-team evaluation against staged
prompt-infection and AgentPoison-class attacks, and measurement of the
practical overhead of the verifier quorum in a cloud-native deployment.
All pseudocode and supporting materials are available at
\url{https://github.com/sunilgentyala/TRACE-MAS}.

%=============================================================================
\begin{thebibliography}{99}

\bibitem{ref:embodied2025}
Z.~Li, W.~Wu, Y.~Guo, J.~Sun, and Q.-L. Han,
``Embodied multi-agent systems: A review,''
\textit{IEEE/CAA J. Autom. Sinica}, vol.~12, no.~6,
pp.~1095--1116, Jun.~2025,
doi:~10.1109/JAS.2025.125552.

\bibitem{ref:tifs2025grey}
Y.~Ma and W.~Li,
``Grey-box adversarial attack on communication in communicative
multi-agent reinforcement learning,''
\textit{IEEE Trans.\ Inf.\ Forensics Security}, vol.~20,
pp.~4679--4693, 2025,
doi:~10.1109/TIFS.2025.3560203.

\bibitem{ref:tifs2024safe}
Z.~Lv, L.~Xiao, Y.~Chen, H.~Chen, and X.~Ji,
``Safe multi-agent reinforcement learning for wireless applications
against adversarial communications,''
\textit{IEEE Trans.\ Inf.\ Forensics Security}, vol.~19,
pp.~6824--6839, 2024,
doi:~10.1109/TIFS.2024.3423428.

\bibitem{ref:access2025multi}
Y.~Hmimou, M.~Tabaa, A.~Khiat, and Z.~Hidila,
``A multi-agent system for cybersecurity threat detection and
correlation using large language models,''
\textit{IEEE Access}, vol.~13, pp.~150199--150215, 2025,
doi:~10.1109/ACCESS.2025.3602681.

\bibitem{ref:jas2025consensus}
C.~Zhou, Z.~Mao, B.~Jiang, and X.-G. Yan,
``Adaptive fault-tolerant consensus tracking control for nonlinear
multi-agent systems with double semi-Markovian switching topologies
and unknown control directions,''
\textit{IEEE/CAA J. Autom. Sinica}, 2025,
doi:~10.1109/JAS.2025.125285.

\bibitem{ref:tc2024pbft}
C.~Li, W.~Qiu, X.~Li, C.~Liu, and Z.~Zheng,
``A dynamic adaptive framework for practical Byzantine fault tolerance
consensus protocol in the Internet of Things,''
\textit{IEEE Trans.\ Comput.}, vol.~73, no.~7,
pp.~1669--1682, Jul.~2024,
doi:~10.1109/TC.2024.3377921.

\bibitem{ref:jas2024priv}
M.~Wei, W.~Yu, D.~Chen, M.~Kang, and G.~Cheng,
``Privacy distributed constrained optimization over time-varying
unbalanced networks and its application in federated learning,''
\textit{IEEE/CAA J. Autom. Sinica}, vol.~12, no.~2,
pp.~335--346, Feb.~2025,
doi:~10.1109/JAS.2024.124869.

\bibitem{ref:tdsc2024}
X.~Zhang, H.~Xu, Z.~Ba, Z.~Wang, Y.~Hong, J.~Liu, Z.~Qin,
and K.~Ren,
``PrivacyAsst: Safeguarding user privacy in tool-using large language
model agents,''
\textit{IEEE Trans.\ Dependable Secure Comput.}, vol.~21, no.~6,
pp.~5242--5258, Nov.--Dec.~2024,
doi:~10.1109/TDSC.2024.3372777.

\bibitem{ref:zpma2024}
A.~Pathak, I.~Al-Anbagi, and H.~J. Hamilton,
``Blockchain-enhanced zero knowledge proof-based privacy-preserving
mutual authentication for IoT networks,''
\textit{IEEE Access}, vol.~12, pp.~118618--118636, 2024,
doi:~10.1109/ACCESS.2024.3450313.

\bibitem{ref:comst2023}
Y.~Wang, T.~Sun, S.~Li, X.~Yuan, W.~Ni, E.~Hossain,
and H.~V. Poor,
``Adversarial attacks and defenses in machine learning-empowered
communication systems and networks: A contemporary survey,''
\textit{IEEE Commun.\ Surveys Tuts.}, vol.~25, no.~4,
pp.~2245--2298, 4th Quart., 2023,
doi:~10.1109/COMST.2023.3319492.

\bibitem{ref:datasentinel2025}
Y.~Liu, Y.~Jia, J.~Jia, D.~Song, and N.~Z. Gong,
``DataSentinel: A game-theoretic detection of prompt injection
attacks,''
in \textit{Proc. 46th IEEE Symp. Security Privacy (SP)},
San Francisco, CA, USA, 2025, pp.~2190--2208,
doi:~10.1109/SP61157.2025.00250.

\bibitem{ref:promptinfection2024}
D.~Lee and M.~Tiwari,
``Prompt infection: LLM-to-LLM prompt injection within multi-agent
systems,'' arXiv preprint arXiv:2410.07283, Oct.~2024. [Online].
Available: \url{https://arxiv.org/abs/2410.07283}

\bibitem{ref:agentpoison2024}
Z.~Chen, Z.~Xiang, C.~Xiao, D.~Song, and B.~Li,
``AgentPoison: Red-teaming LLM agents via poisoning memory or
knowledge bases,''
in \textit{Adv. Neural Inf. Process. Syst. (NeurIPS)}, vol.~37,
2024. [Online]. Available:
\url{https://proceedings.neurips.cc/paper_files/paper/2024/hash/eb113910e9c3f6242541c1652e30dfd6}

\bibitem{ref:agentsec2025}
Y.~Louck, A.~Stulman, and A.~Dvir,
``Security analysis of agentic AI communication protocols:
A comparative evaluation,''
\textit{ACM Trans. AI Security Privacy}, 2025,
doi:~10.1145/3803431.

\bibitem{ref:a2a2025}
Z.~Anbiaee, M.~Rabbani, M.~Mirani, G.~Piya, I.~Opushnyev,
A.~Ghorbani, and S.~Dadkhah,
``Security threat modeling for emerging AI-agent protocols:
A comparative analysis of MCP, A2A, Agora, and ANP,''
arXiv preprint arXiv:2602.11327, Feb.~2026. [Online].
Available: \url{https://arxiv.org/abs/2602.11327}

\end{thebibliography}

\end{document}
